\documentclass{beamer}

% Language and encoding
\usepackage[utf8]{inputenc}
\usepackage[T1]{fontenc}
\usepackage[english]{babel}

% Math packages
\usepackage{amsmath, amssymb, amsthm}
\usepackage{changepage} % Added for adjustwidth environment

% Graphics
\usepackage{graphicx}

% Theme
\usetheme{Madrid}

% Title info
\title{Maps between CW Complexes}
\author{Abraham Rojas Vega}
\date{\today}

% Custom commands (optional)
\newcommand{\R}{\mathbb{R}}
\newcommand{\Z}{\mathbb{Z}}

\begin{document}

\maketitle

% Content goes here

\begin{frame}{CW Complexes}
\begin{enumerate}
    \setlength\itemsep{1.5em}
    \item Start with a discrete set $X^0$, whose points are regarded as 0-cells.
    \item Inductively, form the $n$-skeleton $X^n$ from $X^{n-1}$ by attaching $n$-cells $e_\alpha^n$ via maps $\varphi_\alpha: S^{n-1} \rightarrow X^{n-1}$, denoted by $X^n \cup _f X^{n-1}$ \alert{(Categorical Pushout)}
    \item One can either stop this inductive process at a finite stage, setting $X = X^n$ for some $n < \infty$, or one can continue indefinitely, setting $X = \bigcup_n X^n$ (considering the weak topology).
\end{enumerate}
\end{frame}    

\begin{frame}
“CW” stands for the following two properties shared by any CW-complex: 
\bigskip


\begin{adjustwidth}{1cm}{1cm}

\textbf{C = “closure finiteness”}: a compact subset of a CW-complex intersects the interior of only finitely many cells, hence in particular so does the closure of any cell.
\medskip

\textbf{W = “weak topology”}: A subset of a CW-complex is open or closed precisely if its intersection with the closure of each cell is open or closed, respectively.
\end{adjustwidth}
\bigskip

Whitehead called the interior of the n-disks the “cells”, so that their closure of each cell is the corresponding 
n-disk. \bigskip

\emph{Why do we care about CW complexes?} 

\end{frame}

\begin{frame}{Good properties}
    \vfill
Cell complexes have a very nice mixture of rigidity and flexibility, with enough rigidity to allow many arguments to proceed in a combinatorial cell-by-cell fashion and enough flexibility to allow many natural constructions to be performed on them.
\vfill
\end{frame}

\begin{frame}
    \begin{itemize}
        \setlength\itemsep{1.5em}
        \item Locally contractible, normal paracompact and compactly generated.
        \item Nice categorical properties: quotients, products,
        \item Cellular approximation theorem, that we will see later.
        \item \textbf{Whitehead's theorem: weak equivalence are homotopy equivalences!}
        \item Brown's representability theorem: gives necessary and sufficient conditions for a contravariant functor on the homotopy category of pointed connected CW complexes, to the category of sets, to be representable.
    \end{itemize}
\end{frame}

\begin{frame}{Category of CW-complexes}

A map $f: X \rightarrow Y$ between CW complexes is cellular if $f\left(X^n\right) \subset Y^n$ for all $n \geq 0$.\\
\bigskip
A homotopy $F: X \times I \rightarrow Y$ is cellular if it is a cellular map with respect to the natural CW complex structure on $X \times I$.\\
\bigskip

There are similar definitions for pairs.

\end{frame}

\begin{frame}

\begin{block}{Proposition}
Let $m>n$. Let $U$ be an open subset of $\mathbb{R}^n$ such that the space $\mathrm{cl} U$ is compact. Let $f: \mathrm{cl} U \rightarrow B^m$ be such that $f(\mathrm{fr} U) \subset S^{m-1}$. Then $f$ is homotopic, rel fr $U$, to a map $g$ such that some point of $B^m-S^{m-1}$ is not in the image of $g$.

\end{block}

\bigskip

    \begin{block}{Propostion}
Let $e_\alpha^n$ be an $n$-cell of the $C W$ complex $X$, and let $z \in \overset{\circ}{e^n _\alpha}$. Then $X^n-\overset{\circ}{e^n _\alpha} $ is a strong deformation retract of $X^n-\{z\}$.
    \end{block}
\end{frame}

\begin{frame}{Cellular aproximation theorem}
    \begin{block}{Theorem}
Let $f: X \rightarrow Y$ be a map between CW complexes and let $A$ be a subcomplex of $X$ such that $f \mid A: A \rightarrow Y$ is cellular. Then $f$ is homotopic, rel $A$, to a cellular map.
    \end{block}
\end{frame}

\begin{frame}
\begin{block}{Addendum}
The homotopy $H: X \times I \rightarrow Y$ obtained in the proof has the property that for each cell $e$ of $X, H(e \times I) \subset C(f(e))$.
\end{block}
\bigskip
\begin{block}{Corollary}
    Let $X$ be a CW complex; $X$ is path connected iff $X^1$ is path connected.
\end{block}
\end{frame}

\begin{frame}
\vfill
\centering
\alert{Now, let's talk about categories!}
\vfill
\end{frame}









\end{document}